\section{Teori Permainan dan Model Ising}

% Slide 6.1: Bab 2.3.1 - Pengantar Teori Permainan dan Ekuilibrium Nash
\begin{frame}{Bab 2.3.1 - Teori Permainan}
    \begin{itemize}
        \item Mengkaji bagaimana agen rasional membuat keputusan strategis di bawah ketergantungan agen yang lain.
        \item Terdiri dari beberapa unsur penting yaitu agen, himpunan strategi, dan fungsi utilitas.
        \item \textbf{\textit{Nash Equilibrium:}} Kondisi stabil di mana jika suatu agen mengubah strateginya maka agen tersebut akan mendapatkan utilitas yang lebih rendah.
        \item \textbf{Contoh Game:}
        \begin{enumerate}
            \item Stag Hunt
            \item Prisoner's Dilemma
            \item Battle of the Sexes
        \end{enumerate}
    \end{itemize}
    \begin{figure}
        \centering
        \begin{subfigure}[b]{0.3\textwidth}
            \centering
            \includegraphics[width=\textwidth]{232_payoff_SH.png}
            \caption{\textit{Stag Hunt}}
        \end{subfigure}
        \hfill
        \begin{subfigure}[b]{0.15\textwidth}
            \centering
            \includegraphics[width=\textwidth]{232_payoff_PD.png}
            \caption{\textit{Prisoner's Dilemma}}
        \end{subfigure}
        \hfill
        \begin{subfigure}[b]{0.3\textwidth}
            \centering
            \includegraphics[width=\textwidth]{232_payoff_BoS.png}
            \caption{\textit{Battle of the Sexes}}
        \end{subfigure}
        \caption{Matriks \textit{payoff} klasik (Sumber: Referensi Teori Permainan).}
    \end{figure}
\end{frame}

% Slide 6.2: Bab 2.3.2 - Exact Potential Game (EPG)
\begin{frame}{Bab 2.3.2 - Exact Potential Game (EPG)}
    \begin{itemize}
        \item Merupakan permaianan yang menganggap utilitasnya sebagai fungsi potensial global $\Phi$.
        \item Berfungsi sebagai algoritma umum untuk mencari \textit{Nash Equilibrium} melalui maksimisasi fungsi potensial.
        \item Perubahan utilitas lokal ($\Delta u_i$) merepresentasikan perubahan potensial global ($\Delta \Phi$).
        \item Jika $\Delta \Phi > 0$, maka agen termotivasi untuk melakukan perubahan. Sebaliknya, jika $\Delta \Phi < 0$, maka agen akan menolak untuk melakukan perubahan sehingga didapatkan titik \textit{Nash Equilibrium}.
    \end{itemize}
    \begin{equation}
        \Delta \Phi = \Phi(s_i', s_{-i}) - \Phi(s_i, s_{-i}) = u_i(s_i', s_{-i}) - u_i(s_i, s_{-i})
    \end{equation}
    Dimana,
    \begin{itemize}
        \item $s_i$: Strategi agen $i$
        \item $s_{-i}$: Strategi agen selain agen $i$
        \item $u_i$: Utilitas agen $i$
        \item $u_i'$: Utilitas agen $i$ setelah melakukan perubahan
    \end{itemize}
\end{frame}

% Slide 6.3: Bab 2.3.3 - Penurunan Markowitz ke Kerangka EPG
\begin{frame}{Bab 2.3.3 - Dari Markowitz ke Kerangka EPG}
    \begin{itemize}
        \item \textbf{Transformasi Lagrangian Awal Menjadi Fungsi Potensial Global:}
        \begin{equation*}
            \begin{split}
                \Phi(\omega_i) &= -\frac{1}{\lambda} \mathcal{L}(\omega_i) \\
                &= -\sum_i \mu_i \omega_i + \frac{1}{\lambda} \sum_{i,j} \sigma_{ij} \omega_i \omega_j
            \end{split}
        \end{equation*}
        Dimana,
        \begin{itemize}
            \item $\lambda = \frac{\gamma}{2}$ dimana $\gamma$ merupakan parameter penentu tingkat penghindaran risiko (Risk Aversion) 
            \item $\sum_{i} \omega_i = 1$ yang mana jika ditransformasi ke variabel keputusan $x_i \in \{0, 1\}$, maka $\sum_{i} \omega_i = \frac{\sum_i x_i}{K} = 1$ dimana $K$ merupakan batas jumlah aset yang akan dipilih karena bobot diasumsikan sama rata.
        \end{itemize}
        \item Maka fungsi potensial EPG ($\Phi(x)$) menjadi:
        \begin{equation}
            \Phi(x) = \frac{1}{K} \sum_i \mu_i x_i - \frac{\gamma}{2K^2} \sum_{i,j} \sigma_{ij} x_i x_j
        \end{equation}
    \end{itemize}
\end{frame}

% Slide 7.1: Bab 2.4.1 - Sejarah dan Evolusi Model Ising
\begin{frame}{Bab 2.4.1 - Model Ising}
    \begin{itemize}
        \item \textbf{1920-an:} Wilhelm Lenz dan Ernst Ising merumuskan model interaksi spin diskrit.
        \item \textbf{1944:} Lars Onsager membuktikan transisi fase pada model Ising 2D.
        \item \textbf{Kontemporer:} Ising menjadi "bahasa standar" dalam optimasi kombinatorial karena kemajuan komputasi.
    \end{itemize}
    \begin{center}
        \includegraphics[width=0.5\textwidth]{1D_Model_atT.png} \\
        \footnotesize \textit{Model Ising 1D (Sumber: Figure 2.1 Bab 2).}
    \end{center}
\end{frame}

% Slide 7.2: Bab 2.4.2 - Transformasi QUBO
\begin{frame}{Bab 2.4.2 - Transformasi QUBO}
    \begin{itemize}
        \item Pemetaan fungsi EPG ke operator Hamiltonian: $\hat{H} = -\Phi(x)$
        \item Penurunan bentuk Hamiltonian objek menggunakan identitas $x_i^2 = x_i$
        {\tiny
        \begin{equation}
            \begin{split}
            \hat{H}_{obj}(x) &= \frac{\gamma}{2K^2}\sum_{i=1}^N\sum_{j=1}^N \sigma_{ij} x_i x_j - \sum_{i=1}^N \frac{\mu_i}{K} x_i \\
            &= \frac{\gamma}{2K^2}\left(\sum_{i=1}^N\sigma_i^2 x_i^2 + \sum_{i\ne j}\sigma_{ij} x_ix_j\right) -\sum_{i=1}^N \frac{\mu_i}{K} x_i \\
            &= \frac{\gamma}{2K^2}\sum_{i=1}^N\sigma_i^2 x_i^2 + \frac{\gamma}{2K^2}\sum_{i\ne j}\sigma_{ij} x_ix_j - \sum_{i=1}^N \frac{\mu_i}{K} x_i \\
            &= \frac{\gamma}{2K^2}\sum_{i=1}^N\sigma_i^2 x_i + \frac{\gamma}{2K^2}\sum_{i\ne j}\sigma_{ij} x_ix_j - \sum_{i=1}^N \frac{\mu_i}{K} x_i \\
            &= \sum_{i=1}^N\left(\frac{\gamma \sigma_i^2}{2K^2} - \frac{\mu_i}{K}\right)x_i + \sum_{i\ne j}\frac{\gamma \sigma_{ij}}{2K^2} x_ix_j
            \end{split}
        \end{equation}
        }
    \end{itemize}
\end{frame}

\begin{frame}{Bab 2.4.2 - Transformasi QUBO (Lanjutan)}
    \begin{itemize}
        \item Penambahan suku penalti untuk kendala $\sum x_i = K$:
        {\tiny
        \begin{equation}
            \begin{split}
            \hat{H}_{pen}(x) &= A\left(\sum_{i=1}^N x_i - K\right)^2 \\
            &= A\left(\sum_{i=1}^N x_i^2 + \sum_{i\ne j} x_ix_j - 2K\sum_{i=1}^N x_i + K^2\right) \\
            &= A\sum_{i=1}^N x_i^2 + A\sum_{i\ne j} x_ix_j - 2AK\sum_{i=1}^N x_i + AK^2
            \end{split}
        \end{equation}
        }
        \item Hamiltonian total diperoleh dengan menggabungkan fungsi Hamiltonian objek dengan fungsi penalti:
        {\tiny
        \begin{equation}
            \begin{split}
                \hat{H} =& \hat{H}_{obj} (x) + \hat{H}_{pen} (x) \\
                =& \sum_i \left(\frac{\gamma \sigma_i^2}{2K^2} - \frac{\mu_i}{K}\right)x_i + \sum_{i\ne j}\frac{\gamma \sigma_{ij}}{2K^2} x_ix_j  + A\sum_{i=1}^N x_i^2 + A\sum_{i\ne j} x_ix_j - 2AK\sum_{i=1}^N x_i + AK^2 \\
                =& \sum_i \left(\frac{\gamma \sigma_i^2}{2K^2} + A - \frac{\mu_i}{K} - 2AK\right)x_i + \sum_{i\ne j}\left(\frac{\gamma \sigma_{ij}}{2K^2} + A\right) x_ix_j + AK^2
            \end{split}
        \end{equation}
        }
    \end{itemize}
\end{frame}

\begin{frame}{Bab 2.4.2 - Transformasi QUBO (Lanjutan)}
    \begin{itemize}
        \item Jika dimisalkan
        {\tiny
        \begin{equation*}
            Q_{ii} = \frac{\gamma \sigma_i^2}{2K^2} + A - \frac{\mu_i}{K} - 2AK \quad \text{dan} \quad Q_{ij} = \frac{\gamma \sigma_{ij}}{2K^2} + A
        \end{equation*}
        }
        \item Persamaan QUBO dapat disederhanakan menjadi
        {\tiny
        \begin{equation}
            \hat{H} = \sum_i Q_{ii} x_i + \sum_{i\ne j} Q_{ij} x_ix_j + AK^2
        \end{equation}
        }
        \item Substitusi variabel keputusan $x_i \in \{0, 1\}$ ke spin $\hat{Z}_i \in \{-1, 1\}$
        {\tiny
        \begin{equation}
            \begin{split}
                x_i     &= \frac{1 - \hat{Z}_i}{2} \\
                x_ix_j  &= \frac{1 - \hat{Z}_i - \hat{Z}_j + \hat{Z}_i \hat{Z}_j}{4}
            \end{split}
        \end{equation}
        }
        sehingga menjadi
        {\tiny
        \begin{equation}
            \begin{split}
                \hat{H} &= \sum_i Q_{ii} \left(\frac{1 - \hat{Z}_i}{2}\right) + \sum_{i\ne j} Q_{ij} \left(\frac{1 - \hat{Z}_i - \hat{Z}_j + \hat{Z}_i \hat{Z}_j}{4}\right) + AK^2 \\
                \hat{H} &= \sum_i Q_{ii} \frac{1}{2} - \sum_i Q_{ii} \frac{\hat{Z}_i}{2} + \sum_{i\ne j} Q_{ij} \frac{1}{4} - \sum_{i\ne j} Q_{ij} \frac{\hat{Z}_i}{4} - \sum_{i\ne j} Q_{ij} \frac{\hat{Z}_j}{4} + \sum_{i\ne j} Q_{ij} \frac{\hat{Z}_i \hat{Z}_j}{4} + AK^2 \\
                \hat{H} &= \sum_i \left(- \frac{Q_{ii}}{2} - \sum_{i\ne j} \frac{Q_{ij}}{2}\right) \hat{Z}_i + \sum_{i\ne j} \frac{Q_{ij}}{4} \hat{Z}_i \hat{Z}_j + \left(\sum_i \frac{Q_{ii}}{2} + \sum_{i\ne j} \frac{Q_{ij}}{4} + AK^2\right) \quad \text{karena } Q_{ij} = Q_{ji}
            \end{split}
        \end{equation}
        }
    \end{itemize}
\end{frame}

\begin{frame}{Bab 2.4.2 - Transformasi QUBO (Lanjutan)}
    \begin{itemize}
        \item Bisa kita dapatkan
        \begin{equation*}
            h_i = - \frac{Q_{ii}}{2} - \sum_{i\ne j} \frac{Q_{ij}}{2}
        \end{equation*}
        \begin{equation*}
            J_{ij} = \frac{Q_{ij}}{4}
        \end{equation*}
        \begin{equation*}
            C = \sum_i \frac{Q_{ii}}{2} + \sum_{i\ne j} \frac{Q_{ij}}{4} + AK^2
        \end{equation*}
        \begin{equation*}
            \hat{H} = - \sum_i h_i \hat{Z}_i - \sum_{i<j} J_{ij} \hat{Z}_i \hat{Z}_j + C
        \end{equation*}
        \item \textbf{Dimana;}
        \begin{itemize}
            \item $h_i$ merupakan koefisien interaksi spin dengan fluktuasi pasar.
            \item $J_{ij}$ merupakan koefisien interaksi antar-spin yang bergantung pada kovariansi.
            \item $C$ merupakan konstanta.
        \end{itemize}
    \end{itemize}
\end{frame}
